\begin{frame}{Lecture 3의 Lomuto Partition 재사용}
Fixed-pivot Quickselect는 현재 subarray의 마지막 원소 $A[r]$를 pivot으로 쓴다. partition 뒤
\[A[p..q-1]\le A[q],\qquad A[q+1..r]>A[q].\]
선택한 pivot instance는 최종 index $q$에 고정되고, 현재 subarray에서의 rank는
\[
k=q-p+1
\]
이다. 같은 값의 다른 instance는 왼쪽에 있을 수 있다.
\begin{block}{용어 정리}
Quickselect는 목표가 있는 한쪽만 재귀하는 selection framework다. Fixed-pivot Quickselect는 $A[r]$ 같은 결정적 규칙을 쓰며, 확률 가정 없이는 expected $\Theta(n)$을 보장하지 않는다.
\end{block}
\end{frame}
\begin{frame}{Pivot의 Subarray Rank}
\[k=q-p+1\]
\centering\begin{tikzpicture}[x=.95cm]\foreach \x/\v in {0/8,1/11,2/3}\node[select active]at(\x,0){\v};\node[select pivot]at(3,0){15};\foreach \x/\v in {4/31,5/48,6/20,7/29,8/65,9/73}\node[select active]at(\x,0){\v};\draw[decorate,decoration={brace,mirror,amplitude=5pt}](0,-.65)--(3.4,-.65)node[midway,below=5pt]{$k=4$개};\end{tikzpicture}
\end{frame}
\begin{frame}{세 가지 분기}
\begin{description}\item[$i=k$] pivot $A[q]$가 정답\item[$i<k$] $A[p..q-1]$에서 같은 rank $i$\item[$i>k$] $A[q+1..r]$에서 새 rank $i-k$\end{description}
\selecttakeaway{어느 경우에도 두 recursive subarray를 동시에 탐색하지 않는다.}
\end{frame}
\begin{frame}{QUICKSELECT 의사코드}
\small\begin{algorithmic}[1]
\Require{$p\le r$ and $1\le i\le r-p+1$}
\Procedure{FixedQuickselect}{$A,p,r,i$}
\If{$p=r$}\State \Return $A[p]$\EndIf
\State $q\gets\Call{LomutoPartitionWithPivot}{A,p,r,A[r]}$
\State $k\gets q-p+1$
\If{$i=k$}\State \Return $A[q]$
\ElsIf{$i<k$}\State \Return \Call{FixedQuickselect}{$A,p,q-1,i$}
\Else\State \Return \Call{FixedQuickselect}{$A,q+1,r,i-k$}\EndIf
\EndProcedure
\end{algorithmic}
\end{frame}
\begin{frame}{왜 오른쪽 Rank는 $i-k$인가?}
\centering\begin{tikzpicture}[x=1cm]\node[select discard,minimum width=35mm]at(0,0){left + pivot: $k$개};\node[select target,minimum width=45mm]at(4.5,0){right: 새 rank $i-k$};\draw[-Latex,AlgoOrange,very thick](1.9,0)--(2.3,0);\end{tikzpicture}
\vspace{4mm}오른쪽으로 갈 때는 이미 더 작은 $k$개를 제거했다. 원래 $i$번째는 남은 배열에서 $(i-k)$번째다.
\end{frame}
